\section{One Society, Four Legal Situations}

Between 1970 and 2026 the same body of priests---the Priestly Society of
Saint Pius X, founded by Archbishop Marcel Lefebvre---has stood in four
distinct relations to the authority of the Holy See. It was erected by a
diocesan bishop as a body of the Church's own law, and commended in
writing by a Roman congregation. It was declared suppressed by that
bishop's successor, while its founder denied that the suppression had
ever taken legal effect. It became, after two sets of episcopal
consecrations conferred without a papal mandate, the object of the
gravest penalties in the Church's penal law, and of a formal declaration
that its ministers are in schism. And in between, for a quarter of a
century, it was the object of a Roman policy that mixed penalty with
concession: excommunications remitted, doctrinal discussions opened,
faculties for confession and marriage granted while every act repeated
that the Society itself had no canonical status.

This account asks what the documents actually establish about that
sequence. Who erected the Society, under what law, and with what
juridical effect? What exactly did the acts of 1975 and 1976 decide, and
what does the public record leave unresolved? What did the excommunications
of 1988 attach to, and to whom? What did the remission of 2009 restore,
and what did it explicitly not restore? What is the precise scope of the
faculties granted in 2015, 2016 and 2017? And what did the Holy See
decree, and the Society answer, in July 2026?

\historythesis{The Society's history is a history of acts, and the acts
say less than the polemics on both sides require. Every consequential
turn in this story---the erection of 1 November 1970, the judgment of the
Cardinals' Commission of 6 May 1975 and the suppression that followed,
the suspension \latin{a divinis} of 1976, the protocol of 5 May 1988 and
its repudiation, the excommunications of 1 and 2 July 1988, the remission
of 21 January 2009, the confession and marriage faculties of 2015--2017,
and the decree and explanatory note of 2 July 2026---is documented, and
each document is narrower than the claim usually built on it. Two
questions that both parties treat as settled are, on the public record,
not settled at all: whether the 1975 suppression was procedurally valid,
and whether the recourse lodged on 11 July 2026 suspends the decree of
2~July. This account states what each act says, identifies who said it,
and leaves open what the Holy See has left open.}

The method is documentary. Where an act was printed in \work{Acta
Apostolicae Sedis}, this account reads the fascicle; where the Holy See
publishes an official text or press bulletin, it reads that; where only
the Society's own archive supplies a document, the account says so and
treats the text as a party transcription. Statements of the Holy See are
attributed to the Holy See and statements of the Society to the Society,
in each case with the date and the signatory. No canonical question that
competent authority has left open is decided here, and nothing in this
account is a judgment about the good faith, culpability, or spiritual
state of any person named in it. The evidence classes, terminology, and
the boundaries of the corpus are set out in the terminal appendix.

A word on names. ``Society of Saint Pius X'' (SSPX; French FSSPX;
\latin{Fraternitas Sacerdotalis Sancti Pii X}) is used throughout as the
conventional English designation of the body Archbishop Lefebvre founded;
it carries no implication that the body presently holds, or ever held,
the canonical figure the name might suggest. Roman documents call it
\latin{Fraternitas} or \textit{Fraternità}, usually rendered ``Fraternity''
in the Holy See's own English texts, and quotations preserve whatever form
the cited document uses. ``Traditionalist'' is descriptive, not canonical.
``Valid,'' ``licit,'' ``faculty,'' ``delegation,'' ``censure,'' ``penalty''
and ``status'' are technical terms and are never used as synonyms.
